% questions.tex —— 1958 年·河北卷 · 题目内容片段（无答案）
% 由 tools/generate_projects.py 依来源总账自 1958/1958_河北_数学真题_试题.md 生成，勿手改（改源 MD 后重跑）


\gkpart{一、小题}

\begin{question}
\gkti{1}（本题满分 4 分）求多项式 $x^{376}-2x^{223}+5x^{145}-3x^{28}+3$ 用 $x+1$ 除所得的余数。

\end{question}
\begin{question}
\gkti{2}（本题满分 6 分）试求函数 $y=\sin 2x$ 的周期，并作出它在一个周期内的图象。

\end{question}
\begin{question}
\gkti{3}（本题满分 8 分）求 $\left(\dfrac{1}{\sqrt[3]{x}}+\sqrt{x}\right)^7$ 展开式中含 $x$ 的一次幂的项。

\end{question}
\begin{question}
\gkti{4}（本题满分 12 分）设一圆锥的母线长为 $l$，底面半径为 $r$，用平行于底面的一个平面在这圆锥上截下一个小圆锥，若小圆锥的侧面积等于原来圆锥的侧面积的一半，试求小圆锥的母线的长。

\end{question}
\gkpart{二、大题}

\begin{question}
\gkti{5}（本题满分 16 分）设 $O$ 为 $\triangle ABC$ 内的一点，线段 $AO$ 的延长线交 $BC$ 于 $A_1$ 点，线段 $BO$ 的延长线交 $AC$ 于 $B_1$ 点，线段 $CO$ 的延长线交 $AB$ 于 $C_1$ 点，求证：

（1）$\dfrac{\triangle OBC}{\triangle ABC}=\dfrac{OA_1}{AA_1}$；

（2）$\dfrac{OA_1}{AA_1}+\dfrac{OB_1}{BB_1}+\dfrac{OC_1}{CC_1}=1$。

\end{question}
\begin{question}
\gkti{6}（本题满分 18 分）解方程组：

\nopagebreak[4]
\[
\begin{cases}
\cos 2x+\cos 2y=\dfrac{\sqrt3}{2},\\[2pt]
x+y=\arccos\dfrac12.
\end{cases}
\]

\end{question}
\begin{question}
\gkti{7}（本题满分 18 分）一农场有甲、乙两台打麦机，甲机的工作效率是乙机的 $2$ 倍，先只用甲机打完全部麦子的 $\dfrac23$，然后只用乙机继续打完，所需要的时间比同时用两台打麦机打完全部麦子所需要的时间多 $4$ 天，问分别用一台打麦机打完所有的麦子各需要多少天？

\end{question}
\begin{question}
\gkti{8}（本题满分 18 分）已知方程 $4x^2-(4\cos\theta)x+1=0$ 有实根，求 $\theta$ 的一切可能的值。

\end{question}